\section{Fundamentação Teórica}
\label{sec:fundamentacao}

Modelos modernos de IA generativa, como arquiteturas baseadas em Transformers, utilizam mecanismos de atenção e grandes quantidades de parâmetros para produzir texto, imagens, código ou outros tipos de conteúdo sintético~\cite{vaswani2017attention,brown2020language}. Apesar de esses sistemas serem normalmente discutidos pelo ponto de vista do modelo, uma parte importante da execução aparece em operações de álgebra linear. Multiplicações de matrizes e operações tensoriais estão presentes tanto no treinamento quanto na inferência, o que aproxima essas cargas de problemas clássicos de computação numérica.

A computação de alto desempenho busca reduzir o tempo de execução de problemas custosos por meio do uso eficiente dos recursos disponíveis. Nesse tipo de análise, não basta observar apenas o nome do processador ou a quantidade de memória instalada. O desempenho também depende da hierarquia de memória, do paralelismo explorado, da movimentação dos dados e das bibliotecas usadas pela aplicação.

No caso de IA generativa, o treinamento é influenciado pela quantidade de parâmetros, pelo tamanho do conjunto de dados, pela precisão numérica e pela estratégia de paralelização. Na inferência, tornam-se relevantes a latência de resposta, a vazão de requisições, o consumo de memória para armazenar pesos do modelo e a movimentação dos dados durante a geração. Por isso, mesmo sem executar um modelo completo, uma carga matricial permite observar uma parte concreta desse comportamento em escala controlada.

Bibliotecas de álgebra linear, como BLAS, são uma base importante para esse tipo de medição, pois padronizam rotinas comuns de álgebra linear, incluindo operações vetor-vetor, matriz-vetor e matriz-matriz~\cite{blackford2002blas,netlibBlas}. Na prática, bibliotecas numéricas usadas em Python, como NumPy, podem se apoiar em implementações otimizadas dessas rotinas para obter desempenho superior ao de laços escritos diretamente em linguagem interpretada.

Ambientes de HPC e IA têm se aproximado porque ambos dependem de processamento paralelo, aceleradores, bibliotecas otimizadas e infraestrutura capaz de reduzir o tempo até a obtenção de resultados~\cite{huerta2020convergence}. Para este trabalho, essa aproximação é importante porque desloca a discussão de IA para além do algoritmo: a execução também depende de CPU, GPU, memória, armazenamento, energia e refrigeração funcionando como partes de um mesmo fluxo computacional.

Benchmarks e medições técnicas ajudam a transformar essa discussão em dados observáveis. Métricas como tempo de execução, taxa estimada de operações por segundo e variação entre repetições permitem interpretar o comportamento de uma carga computacional sem depender apenas de especificações do fabricante. No experimento deste artigo, essas métricas foram usadas para avaliar uma carga matricial em CPU e comparar sua escala com a de sistemas dedicados.

Como referência de infraestrutura especializada, o sistema NVIDIA DGX H100 é documentado com 8 GPUs H100, 640 GB de memória total de GPU, 2 TB de memória do sistema, desempenho de 32 petaFLOPS em FP8 e consumo máximo aproximado de 10,2 kW~\cite{nvidiaDgxH100}. Esses valores não representam o ambiente experimental deste trabalho, mas ilustram a escala de recursos normalmente associada a cargas modernas de IA em infraestrutura de alto desempenho.
